%! Observational Studies — Unit 1, Lesson 1.5 — Warm-Up (Answer Key)
% ONE PAGE, blank and key. Three items at 12pt (the stats-model size), \workrowsep 50pt, item
% spacing 22pt, so each item gets real handwriting room and still fits. Do not add a fourth
% item — the page is full.
% GRADUAL RELEASE: the warm-up activates prior knowledge before the guided notes. All three
% items sit on the Harbor Point sample the notes open with, so 70%, 45%, 25 points, and 55%
% are on the board before the word "observational" is spoken. Items 1 and 2 are the two group
% rates WITH TWO DIFFERENT DENOMINATORS (80 and 120) — row 3 of the notes (problem 9) reads
% them straight off the two-way table; item 3 is the gap and the pooled rate (problem 10 scales
% the 55% to all 1,500). The tell is item 3's pooled rate: (70 + 45)/2 = 57.5% merges two
% unequal groups.
\documentclass[12pt]{article}
\usepackage{saar-article}
\usepackage{saar-key}   % pulls in -boxes; do NOT also load -boxes
\newcommand{\TallMath}[1]{$\displaystyle #1\rule[-1.4em]{0pt}{3.2em}$}
% Word bank strip for every fill-in cluster (user request 2026-09-06). Defined per document;
% the key inherits it and prints the same strip, so heights match.
\newcommand{\wordbank}[1]{\par\vspace{2pt}\noindent\fcolorbox{goldacc}{goldbg}{\parbox{\dimexpr\linewidth-2\fboxsep-2\fboxrule\relax}{\small\textbf{Word bank:} #1}}\par\vspace{4pt}}
% Handwriting room inside every \begin{work} block. Moves the blank and the
% key together, so the two can never drift apart.
\setlength{\workrowsep}{50pt}

\begin{document}

\pageheader{Unit 1, Lesson 1.5}{Warm-Up --- Answer Key}
% NAMESTRIP: no \namedateperiod here — mirrors the blank. NO teachernote here —
% teacher prose lives in the lesson plan.

\vspace{0.15in}

\begin{notesbox}{Spiral Review --- Back to the Harbor Point Community Pool}
\normalsize
Harbor Point drew a random sample of $200$ of its $1{,}500$ members from the full membership
list, and every one of them answered. Of the $200$, $80$ swim in the morning and $120$ swim
in the evening. Each was asked whether they usually get at least $7$ hours of sleep.

\begin{enumerate}[leftmargin=1.8em, itemsep=22pt, topsep=10pt]

  \item $56$ of the $80$ morning swimmers said yes. What \textbf{percent} of the morning
        swimmers is that?

        \begin{work}
          \dfrac{56}{80} &= 0.70 = 70\%
        \end{work}

  \item $54$ of the $120$ evening swimmers said yes. What \textbf{percent} of the evening
        swimmers is that?

        \begin{work}
          \dfrac{54}{120} &= 0.45 = 45\%
        \end{work}

  \item Compare the two groups, then put the whole sample back together.
        \wordbank{25 \;\textbullet\; 110 \;\textbullet\; 55\% \;\textbullet\; 57.5\% \;\textbullet\; 80 \;\textbullet\; 120 \;\textbullet\; 200}

        \vspace{4pt}
        \renewcommand{\arraystretch}{2.6}
        \begin{tabularx}{\linewidth}{@{} X p{3.4cm} @{}}
        \toprule
        \textbf{Question} & \textbf{Answer} \\
        \midrule
        How many \textbf{percentage points} higher is the morning group? & \ans{25} \\
        How many of the $200$ members in the sample said yes in total?  & \ans{110} \\
        What \textbf{percent of the whole sample} of $200$ said yes?     & \ans{$110/200 = 55\%$} \\
        \bottomrule
        \end{tabularx}

\end{enumerate}
\end{notesbox}

\end{document}
